\begin{activity}{Activity 1.5: Deciding Whether the Binomial Theorem Can Be Used in All Cases}
Discuss in groups whether the binomial theorem
\[
  (1+x)^n = 1 + nx + \frac{n(n-1)}{2!}x^2 + \frac{n(n-1)(n-2)}{3!}x^3 + \cdots + x^n
\]
can be used in all cases.
\begin{enumerate}
  \item Consider $(2+4x)^4$. Can the binomial theorem be used directly?
  \item How will you write the terms in the bracket to have $1$ as one of the terms?
  \item Let us go through the following tasks to help us write $(2+4x)^4$ in the
  form $(1+x)^n$ and use the binomial theorem to expand and simplify:
  \begin{itemize}
    \item[] \textbf{Task 1:} Take the given binomial expression $(2+4x)$.
    \item[] \textbf{Task 2:} Factorise the constant out from the binomial to obtain $2(1+2x)$.
    \item[] \textbf{Task 3:} Rewrite your expression obtained with the power as given, i.e.\ $2^4(1+2x)^4$.
    \item[] \textbf{Task 4:} Expand the binomial $(1+2x)^4$, after which you multiply your result by $2^4$.
  \end{itemize}
  Now follow the tasks to obtain all the terms for the given expression.

  \textbf{Expected Answer:} $16 + 128x + 348x^2 + 512x^3 + 256x^4$.
\end{enumerate}
\end{activity}

Let's solve the following examples using the binomial theorem.

\begin{example}{Example 1.16}
\begin{enumerate}
  \item Use the binomial theorem to expand fully $(3+9x)^5$.
  \item Find the first three terms of $\left(\dfrac{4}{9}+6x\right)^{1/2}$ using the binomial theorem.
  \item Use the binomial theorem to expand fully and simplify $(10+4x)^3$.
  \item Use the binomial theorem to find the first four terms of the expression $\sqrt[4]{81+162x}$.
\end{enumerate}

\medskip\noindent\textbf{Solution}

\textbf{1.}
\[
  (3+9x)^5 = 3^5(1+3x)^5
\]
\begin{align*}
  3^5(1+3x)^5 &= 3^5\Bigg[1 + 5(3x) + \frac{5(4)}{2\times1}(3x)^2 + \frac{5(4)(3)}{3\times2\times1}(3x)^3 \\
  &\qquad + \frac{5(4)(3)(2)}{4\times3\times2\times1}(3x)^4 + \frac{5(4)(3)(2)(1)}{5\times4\times3\times2\times1}(3x)^5\Bigg]
\end{align*}
\[
  = 243\left[1 + 5(3x) + 10(9x^2) + 10(27x^3) + 5(81x^4) + (243x^5)\right]
\]
\[
  (3+9x)^5 = 243\left(1 + 15x + 90x^2 + 270x^3 + 405x^4 + 243x^5\right)
\]
Compute the final coefficients:
\begin{itemize}
  \item $243 \times 1 = 243$
  \item $243 \times 15 = 3645$
  \item $243 \times 90 = 21870$
  \item $243 \times 270 = 65610$
  \item $243 \times 405 = 98415$
  \item $243 \times 243 = 59049$
\end{itemize}
Therefore, the expansion is
\[
  (3+9x)^5 = 243 + 3645x + 21870x^2 + 65610x^3 + 98415x^4 + 59049x^5.
\]

\textbf{2.}
\[
  \left(\frac{4}{9}+6x\right)^{1/2} = \left(\frac{4}{9}\right)^{1/2}\left(1+\frac{27}{2}x\right)^{1/2}
\]
\[
  = \left(\frac{4}{9}\right)^{1/2}\left[1 + \frac{1}{2}\left(\frac{27}{2}x\right) + \frac{\tfrac{1}{2}\left(\tfrac{1}{2}-1\right)}{2!}\left(\frac{27}{2}x\right)^2\right]
\]
\[
  = \frac{2}{3}\left[1 + \frac{27}{4}x + \frac{-\tfrac{1}{4}}{2}\left(\frac{729}{4}x^2\right)\right]
  = \frac{2}{3}\left[1 + \frac{27}{4}x - \frac{729}{32}x^2\right]
\]
\[
  = \frac{2}{3} + \frac{9}{2}x - \frac{243}{16}x^2
\]

\textbf{3.}
\[
  (10+4x)^3 = 10^3\left(1+\frac{2}{5}x\right)^3
\]
\[
  = 10^3\left[1 + 3\left(\frac{2}{5}x\right) + \frac{3(2)}{2\times1}\left(\frac{2}{5}x\right)^2 + \frac{3(2)(1)}{3\times2\times1}\left(\frac{2}{5}x\right)^3\right]
\]
\[
  = 10^3\left[1 + \frac{6}{5}x + 3\left(\frac{4x^2}{25}\right) + \frac{8}{125}x^3\right]
  = 1000\left[1 + \frac{6}{5}x + \frac{12x^2}{25} + \frac{8x^3}{125}\right]
\]
\[
  (10+4x)^3 = 1000 + 1200x + 480x^2 + 64x^3
\]

\textbf{4.}
\[
  \sqrt[4]{81+162x} = (81+162x)^{1/4} = 81^{1/4}(1+2x)^{1/4}
\]
\[
  = 81^{1/4}\left[1 + \frac{1}{4}(2x) + \frac{\tfrac{1}{4}\left(\tfrac{1}{4}-1\right)}{2!}(2x)^2 + \frac{\tfrac{1}{4}\left(\tfrac{1}{4}-1\right)\left(\tfrac{1}{4}-2\right)}{3!}(2x)^3\right]
\]
\[
  = 81^{1/4}\left[1 + \frac{1}{4}(2x) + \frac{-\tfrac{3}{16}}{2\times1}(4x^2) + \frac{\tfrac{21}{64}}{3\times2\times1}(8x^3)\right]
  = 3\left[1 + \frac{1}{2}x - \frac{3}{8}x^2 + \frac{7}{16}x^3\right]
\]
\[
  \sqrt[4]{81+162x} = 3 + \frac{3}{2}x - \frac{9}{8}x^2 - \frac{21}{16}x^3
\]
\end{example}

\begin{example}{Example 1.17}
The first three terms in the expansion of $\left(1+\dfrac{x}{p}\right)^n$ in ascending powers
of $x$ are $1 + x + \dfrac{9}{20}x^2$. Find the values of $n$ and $p$.

\medskip\noindent\textbf{Solution}

Expand $\left(1+\dfrac{x}{p}\right)^n$ using the binomial theorem:
\[
  \left(1+\frac{x}{p}\right)^n = 1 + n\frac{x}{p} + \frac{n(n-1)}{2!}\left(\frac{x}{p}\right)^2 = 1 + nx + \frac{n(n-1)}{2}\cdot\frac{x^2}{p^2}
\]
Equating this to $1 + x + \dfrac{9}{20}x^2$ and comparing the coefficients of $x$ and $x^2$:
\[
  \frac{n}{p} = 1, \qquad \frac{n(n-1)}{2p^2} = \frac{9}{20}
\]
\[
  n = p \quad \ldots (1) \qquad\qquad 20n(n-1) = 18p^2 \quad \ldots (2)
\]
Substituting $n=p$ from (1) into (2):
\[
  20p(p-1) = 18p^2 \implies 20p^2 - 20p = 18p^2 \implies 2p^2 = 20p \implies p = 10
\]
\[
  \therefore\ p = 10, \qquad n = 10
\]
\end{example}

\begin{example}{Example 1.18}
Find the values of $t$ and $n$ in the equation $(1+tx)^n = 1 - 6x + \dfrac{33}{2}x^2$.

\medskip\noindent\textbf{Solution}

Expand $(1+tx)^n$ using the binomial theorem:
\[
  (1+tx)^n = 1 + ntx + \frac{n(n-1)}{2!}(tx)^2 = 1 + ntx + \frac{n(n-1)}{2}t^2x^2
\]
Equating this to $1 - 6x + \dfrac{33}{2}x^2$ and comparing coefficients of $x$ and $x^2$:
\[
  nt = -6 \quad \ldots (1) \qquad\qquad \frac{n(n-1)}{2}t^2 = \frac{33}{2} \implies 2n(n-1)t^2 = 66 \quad \ldots (2)
\]
Rewriting (2): $(2n^2-2n)t^2 = 66 \implies 2(nt)^2 - 2(nt)t = 66$.
Since $nt=-6$:
\[
  2(-6)^2 - 2(-6)t = 66 \implies 72 + 12t = 66 \implies 12t = -6 \implies t = -\frac{1}{2}
\]
Substituting back into $nt=-6$:
\[
  n\left(-\frac{1}{2}\right) = -6 \implies -n = -12 \implies n = 12
\]
\end{example}

\begin{example}{Example 1.19}
Expand $(1+2x-x^2)^6$ as far as the $x^2$ term.

\medskip\noindent\textbf{Solution}

How will you use the binomial theorem to solve a question of this nature?
Discuss your views with a classmate or in your groups.

You need to write the trinomial $1+2x-x^2$ as a binomial in the form $(1+x)^n$.
We represent $2x - x^2$ by a variable, say $y$, i.e.\ $y = 2x - x^2$, so
\[
  (1+2x-x^2)^6 = (1+y)^6.
\]
Expanding $(1+y)^6$ using the binomial theorem:
\[
  (1+y)^6 = 1 + 6y + \frac{6(6-1)}{2!}y^2 + \cdots = 1 + 6y + \frac{6(5)}{2}y^2 = 1 + 6y + 15y^2
\]
But $y = 2x-x^2$, so substituting:
\[
  = 1 + 6(2x-x^2) + 15(2x-x^2)^2 + \cdots
\]
\[
  = 1 + 6(2x-x^2) + 15(4x^2 - 4x^3 + x^4) + \cdots
\]
\[
  = 1 + 12x - 6x^2 + 60x^2 - 60x^3 + 15x^4 + \cdots
\]
\[
  (1+2x-x^2)^6 = 1 + 12x + 54x^2 + \cdots
\]
\end{example}

\begin{example}{Example 1.20}
Expand $(1-3x+x^2)^5$ as far as the $x^3$ term.

\medskip\noindent\textbf{Solution}

We represent $-3x+x^2$ by a variable, say $y$, i.e.\ $y = -3x+x^2$, so
\[
  (1-3x+x^2)^5 = (1+y)^5.
\]
Expanding $(1+y)^5$ using the binomial theorem:
\[
  (1+y)^5 = 1 + 5y + \frac{5(4)}{2}y^2 + \frac{5(4)(3)}{6}y^3 = 1 + 5y + 10y^2 + 10y^3
\]
But $y = -3x+x^2$, so substituting:
\[
  = 1 + 5(-3x+x^2) + 10(-3x+x^2)^2 + 10(-3x+x^2)^3
\]
\[
  = 1 + 5(-3x+x^2) + 10(9x^2 - 6x^3 + x^4) + 10(-27x^3 + 27x^4 - 9x^5 + x^6)
\]
\[
  = 1 - 15x + 5x^2 + 90x^2 - 60x^3 + 10x^4 - 270x^3 + 270x^4 - 90x^5 + 10x^6
\]
Collecting terms up to $x^3$:
\[
  (1-3x+x^2)^5 = 1 - 15x + 95x^2 - 330x^3 + \cdots
\]
\end{example}

Did you know we can use binomial expansion to approximate exponential
numbers? This powerful technique allows us to simplify complex expressions
and make calculations easier. We will explore how to apply binomial expansion
to approximate exponential functions.
